\documentclass[aps,prd,preprintnumbers,nofootinbibn,twocolumn]{revtex4}
%showpacs,superscriptaddress,groupedaddress
\usepackage{graphicx}  % needed for figures
\usepackage{dcolumn}   % needed for some tables
\usepackage{bm}        % for math
\usepackage{amssymb}   % for math
\usepackage[utf8]{inputenc}
%\usepackage[spanish]{babel}
\usepackage{hyperref}

\usepackage{amsmath}
\usepackage{array}
\usepackage{appendix}
\usepackage{epsfig}
\usepackage{siunitx}

\hyphenation{arXiv}

\addtolength{\textheight}{1.6cm}

\newcommand{\dif}{\mathrm{d}}
\makeatletter
\def\l@subsection#1#2{}
\def\l@subsubsection#1#2{}
\makeatother
\begin{document}

%%Nuevos desarrollos en teorías de campos no locales
%\title{Fenomenología cosmológica de cuerdas p-ádicas}
%\input autores.tex
%\date{\today}


%\begin{abstract}
%En esta memoria se analiza el modelo de cuerdas p-ádicas a temperatura finita desde un punto de vista cosmológico. En %particular, se estudia la evolución de sus propiedades termodinámicas en un universo en expansión. La dependencia de la %temperatura del fluido p-ádico con el factor de escala del universo y su contenido de energía es fundamental para la comprensión del fluido p-ádico como el origen de la inflación en el universo más temprano o como energía oscura en la evolución cosmológica tardía. Se presta especial atención a la realización particular de esta posibilidad  y el efecto sobre el fondo cósmico de microondas.

%The results presented here have been obtained based on the theory of  finite temperature for p-adic string models. First we found the exact expressions of the pressure, energy density  and entropy density of a p-adic string fluid, for which has been obtained the parameter of the state equation too. Secondly, we studied the temperature of the fluid in terms of the scale factor of the universe, emphasizing the high and low temperature limits. Finally, some  notes relating the p-adic string gas temperature fluctuations with cosmic microwave background are given: in order to find how thermal fluctuations in a cosmological context define the power spectrum  ${\cal P}_{\Phi}(k)$, it had to be evaluated at the instant when the mode $k$ crossed the  Hubble horizon. A formula for $k(T)$ is calculated for a universe dominated by the p-adic fluid, by a cosmological constant, by matter and by radiation.
%\end{abstract}

%\pacs{}
%\maketitle


\title{NEW APPLICATIONS OF THE COLEMAN-WEINBERG MODEL\\[0.4cm]}
\thispagestyle{empty}


\date{de 2016}


\begin{abstract}
\begin{center}
\vspace{0.7cm}

{\bf Autor:}\\

\vspace{0.2cm}

Jorge Alda Gallo $^{(1)}$\\

\vspace{1cm}

{\bf Director:}\\

\vspace{0.2cm}

J. A. R. Cembranos $^{(2)}$\\

\vspace{0.6cm}


\vspace{5cm}
\begin{center}
\includegraphics[width=0.30\textwidth]{ucmlogo}
\end{center}
\vspace{5cm}
\end{center}

$^{(1)}$  E-mail: \href{mailto:jalda@ucm.es}{jalda@ucm.es}

$^{(2)}$ E-mail: \href{mailto:cembra@fis.ucm.es}{cembra@fis.ucm.es}
\vspace{8cm}

\end{abstract}




\clearpage


\maketitle

\begin{widetext}
\begin{quotation}

{\bf Resumen:}\\



\vspace{2cm}

{\bf Abstract:}\\



\vspace{2cm}
\end{quotation}
\end{widetext}
\tableofcontents


\vspace{5cm}




\newpage
\clearpage

\section{Introduction}
The Higgs boson has remained for a long time the elusive last piece of the standard model of particle physics (SM) to be discovered. The waiting came to an end on 2012, when it was announced the discovery at LHC \cite{Aad:2012tfa,Chatrchyan:2012xdj} of a new resonance at \SI{125}{\giga\electronvolt}.

But a \SI{125}{\giga\electronvolt} SM Higgs boson is not without problems. The coupling of the top quark with the Higgs might mean that the electroweak vacuum is metastable \cite{EliasMiro:2011aa}, and it would be able to decay to a lower energy vacuum. Naturalness indicates that the Higgs mass is far too low compared to the Planck scale, which leads to the so-called hierarchy problem \cite{Iso:2013aqa}. 


According to the SM, the Higgs field $H$ is characterized by the following potential
\begin{equation}
V(H) = m^2 H^\dagger H + \lambda_h (H^\dagger H)^2
\end{equation}
Here, $m^2 < 0$, which means that $H=0$ is unstable. Thus, the Higgs field has a vacuum expectation (vev) $v=\SI{246}{\giga\electronvolt}$ (known from masses of the $W$ and $Z$ bosons). The condition for the minimum of the potential is 
\begin{equation}
-m^2 = \lambda_h v^2
\end{equation}
The physical mass of the Higgs boson is determined as the value of the second derivative of the potential at the vacuum, that is, 
\begin{equation}
m_h^2 = 2 \lambda_h v^2
\end{equation}

An alternative proposal for the origin of the Higgs boson is due to Coleman and Weinberg \cite{Coleman:1973jx}. The mass term is set to zero, and the symmetry breaking is caused by radiative corrections to the quartic term. In this way, the Coleman-Weinberg lagrangian is classically scale invariant and is free of the hierarchy problem.

\section{The model} 

\bibliographystyle{apsrev4-1}
\bibliography{memoria} 

\newpage
\clearpage

\begin{widetext}
El abajo firmante, matriculado en el Máster de Física Teórica de la Facultad de Ciencias Físicas, autoriza a la Universidad Complutense de Madrid (UCM) a difundir y utilizar con fines académicos, no comerciales y mencionando expresamente a su autor el presente Trabajo de Fin de Máster: Nuevas aplicaciones del modelo de Coleman-Weinberg, realizado durante el curso académico 2015-2016 bajo la dirección de Jose A. Ruiz Cembranos en el Departamento de Física Teórica I, y a la Biblioteca de la UCM a depositarla en el Archivo institucional E-Prints Complutense u otra plataforma de la UCM que se cree para tal fin con el objeto de incrementar la difusión, uso e impacto del trabajo en Internet y garantizar su preservación y acceso a largo plazo.\\

La publicación en abierto tendrá un embargo de:
\\

\noindent$\,X\,$ Ninguno
\\


Fdo:\\


\vspace{3cm}

El abajo firmante, director del Máster de Física Teórica de la Facultad de Ciencias Físicas, autoriza a la Universidad Complutense de Madrid (UCM) a difundir y utilizar con fines académicos, no comerciales y mencionando expresamente a su autor el presente Trabajo de Fin de Máster: Nuevas aplicaciones del modelo de Coleman-Weinberg, realizado durante el curso académico 2015-2016 bajo mi dirección en el Departamento de Física Teórica I, y a la Biblioteca de la UCM a depositarla en el Archivo institucional E-Prints Complutense u otra plataforma de la UCM que se cree para tal fin con el objeto de incrementar la difusión, uso e impacto del trabajo en Internet y garantizar su preservación y acceso a largo plazo.\\

La publicación en abierto tendrá un embargo de:
\\

\noindent$\,X\,$ Ninguno
\\


Fdo:\\


\vspace{12cm}

\end{widetext}



\end{document} 
